% =============================================================================
%  STUDY COMPANION — Foundations of Large Language Models
%  ZG536 LLMs for GenAI, Session 1
%  Built with make-lecture-kit v3.1.0
% =============================================================================

\documentclass[a4paper,11pt]{article}

\usepackage[T1]{fontenc}
\usepackage[utf8]{inputenc}
\IfFileExists{mathpazo.sty}{\usepackage[sc]{mathpazo}\linespread{1.06}}{\usepackage{lmodern}}
\IfFileExists{helvet.sty}{\usepackage[scaled=0.9]{helvet}}{}
\renewcommand{\ttdefault}{lmtt}
\usepackage{microtype}

\usepackage[a4paper,top=26mm,bottom=26mm,left=27mm,right=27mm,
            headheight=14pt,headsep=12mm,footskip=12mm]{geometry}

\usepackage{amsmath,amssymb}

\usepackage{graphicx}
\usepackage{booktabs}
\usepackage{array}
\usepackage[table,svgnames]{xcolor}
\usepackage{caption}
\usepackage{float}
\usepackage[section]{placeins}
\usepackage{needspace}

\usepackage[most]{tcolorbox}
\tcbuselibrary{skins,breakable}

\usepackage{tikz}
\usepackage{pifont}

\usepackage{enumitem}
\setlist{leftmargin=1.5em,topsep=4pt,itemsep=3pt}

\newlist{steps}{enumerate}{1}
\setlist[steps]{label={\sffamily\small\bfseries\color{accent}Step \arabic*},
  leftmargin=4.2em, labelwidth=3.6em, labelsep=0.6em, labelindent=0pt,
  align=left, topsep=6pt, itemsep=5pt}
\usepackage{parskip}
\setlength{\parskip}{0.55em plus 0.1em}

\usepackage{fancyhdr}
\usepackage{titlesec}

\usepackage[hidelinks]{hyperref}
\usepackage{xurl}
\hypersetup{breaklinks=true}

\definecolor{accent}{HTML}{1F4E79}
\definecolor{inktext}{HTML}{1F2328}
\definecolor{mutedtext}{HTML}{6B7280}
\definecolor{hairline}{HTML}{D0D5DD}
\definecolor{wash}{HTML}{F5F6F8}
\colorlet{navy}{accent}      \colorlet{navyrule}{hairline}
\colorlet{intuFrame}{accent} \colorlet{everyFrame}{accent} \colorlet{workFrame}{accent}
\colorlet{watchFrame}{inktext} \colorlet{keyFrame}{accent} \colorlet{waitFrame}{accent}

\color{inktext}
\hypersetup{linkcolor=accent,urlcolor=accent,citecolor=accent}

\newcommand{\SeriesName}{ZG536 LLMs for GenAI}
\newcommand{\TopicTitle}{Foundations of Large Language Models}

\pagestyle{fancy}
\fancyhf{}
\fancyhead[L]{\sffamily\scriptsize\color{mutedtext}\MakeUppercase{\SeriesName}}
\fancyhead[R]{\sffamily\scriptsize\color{mutedtext}\TopicTitle}
\fancyfoot[C]{\sffamily\footnotesize\color{mutedtext}\thepage}
\renewcommand{\headrulewidth}{0pt}
\renewcommand{\footrulewidth}{0pt}
\fancypagestyle{plain}{\fancyhf{}\fancyfoot[C]{\sffamily\footnotesize\color{mutedtext}\thepage}}

\titleformat{\section}
  {\sffamily\Large\bfseries\color{inktext}}
  {\color{accent}\thesection}{0.7em}{}
\titleformat{\subsection}
  {\sffamily\large\bfseries\color{inktext}}
  {\color{accent}\thesubsection}{0.6em}{}
\titlespacing*{\section}{0pt}{26pt}{6pt}
\titlespacing*{\subsection}{0pt}{14pt}{4pt}
\newcommand{\sectionbreak}{\needspace{14\baselineskip}}

\newcommand{\secsub}[1]{%
  \noindent{\large\itshape\color{mutedtext}#1}\par
  \vspace{2pt}{\color{hairline}\rule{\linewidth}{0.4pt}}\par\medskip}

\newcommand{\flipanswer}[1]{\par\noindent\hfill\rotatebox[origin=c]{180}{\parbox{0.9\linewidth}{\small\color{mutedtext}#1}}\par\medskip}

\tcbset{
  quiet/.style 2 args={
    enhanced, breakable, sharp corners,
    colback=#2, colframe=#2, boxrule=0pt,
    borderline west={1.6pt}{0pt}{#1},
    left=12pt, right=10pt, top=7pt, bottom=7pt,
    before skip=12pt, after skip=12pt,
    fontupper=\normalsize, coltitle=#1,
    fonttitle=\sffamily\scriptsize\bfseries,
    attach title to upper={\par\smallskip},
  }
}
\newcommand{\calloutlabel}[1]{\MakeUppercase{#1}}

\newtcolorbox{intuition}{quiet={accent}{white},
  title={\calloutlabel{The intuition}}}
\newtcolorbox{everyday}{quiet={accent}{white},
  title={\calloutlabel{Everyday picture}}}
\newtcolorbox{worked}[1]{quiet={accent}{wash},
  title={\calloutlabel{Worked example}\ \textperiodcentered\ {\normalfont\itshape\small\color{inktext}#1}}}
\newtcolorbox{watchout}{quiet={inktext}{white},
  title={\calloutlabel{Watch out}}}
\newtcolorbox{waitwhat}{quiet={accent}{white},
  title={\calloutlabel{Wait --- really?}}}
\newtcolorbox{keytake}{quiet={accent}{white},
  borderline north={0.4pt}{0pt}{hairline},
  borderline south={0.4pt}{0pt}{hairline},
  title={\upshape\calloutlabel{Key takeaway}}, fontupper=\normalsize\itshape}

\newcommand{\titlebanner}[4]{%
  \begingroup\raggedright
  \vspace*{4mm}
  {\sffamily\footnotesize\bfseries\color{accent}\MakeUppercase{#1}}\par
  \vspace{7pt}
  {\sffamily\bfseries\fontsize{30}{35}\selectfont\color{inktext}#2\par}
  \vspace{7pt}
  {\large\color{mutedtext}#3\par}
  \vspace{12pt}
  {\color{hairline}\rule{\linewidth}{0.5pt}}\par
  \vspace{3pt}
  {\sffamily\scriptsize\color{mutedtext}\MakeUppercase{#4}}\par
  \vspace{10pt}
  \endgroup
}

\captionsetup[figure]{font={small,color=mutedtext},labelfont={bf,sf,color=inktext},
  labelsep=period,justification=raggedright,singlelinecheck=false,skip=6pt}
\setlength{\textfloatsep}{16pt plus 4pt minus 4pt}
\setlength{\intextsep}{14pt plus 4pt minus 4pt}
\renewcommand{\topfraction}{0.8}\renewcommand{\bottomfraction}{0.6}
\renewcommand{\textfraction}{0.15}\renewcommand{\floatpagefraction}{0.75}
\newcommand{\housefig}[3][0.86]{%
  \begin{figure}[!htbp]
    \centering
    \includegraphics[width=#1\linewidth]{#2}%
    \caption{#3}
  \end{figure}%
}

\renewcommand{\arraystretch}{1.18}
\newcommand{\vocab}[1]{\textbf{#1}}

\begin{document}

\thispagestyle{fancy}

\titlebanner
  {\SeriesName}
  {Foundations of Large Language Models}
  {From a plain sentence to a predicted next word: the five-stage pipeline inside every LLM}
  {Built from: the Session 1 handout, slides and recording · Every example worked out in full}

\noindent\textbf{How to use this companion.}
This companion explains Session 1 from scratch, in plain English. Every idea
starts with an everyday picture. Then the mechanism is built up one step at a
time. Every example is worked out in full, with real numbers.
\textsf{\small\textbf{THE INTUITION}} states the core idea in one breath.
\textsf{\small\textbf{EVERYDAY PICTURE}} ties it to real life before any
mechanism. \textsf{\small\textbf{WORKED EXAMPLE}}, on a grey background,
walks a real calculation through with no skipped steps. \textsf{\small\textbf{WATCH
OUT}} flags the common mistake. \textsf{\small\textbf{WAIT --- REALLY?}}
holds a result that sounds wrong but is true. \textsf{\small\textbf{KEY
TAKEAWAY}} is the one sentence to remember. Read it in order; each section
leans on the one before.

\medskip

% =============================================================================
\section{Try it first: trace one sentence through an LLM}
\secsub{Before any theory, follow one short sentence as far as you can, by
feel, and see exactly where the trail runs out.}

Take the sentence \textit{``The dog loved the movie''} and imagine handing it
to a large language model. Underneath everything, that model is just a very
large pile of numbers and arithmetic. Three honest questions:

\begin{quote}
\itshape
How does a model built only from numbers ``read'' the word \textit{dog}? What
does it mean for the model to ``remember'' that \textit{dog} came before
\textit{movie}, not after? And once it has read the whole sentence, how does
it pick the next word?
\end{quote}

You cannot answer any of these yet, and that is fine --- by the end of this
session you can. Every LLM, from the smallest to the largest, pushes raw text
through the same five-stage pipeline, once per word it generates.

\housefig{figures/fig_pipeline.png}{The five-stage pipeline every sentence
travels through. Sections 2 to 5 build each stage, one at a time.}

Section 2 explains what ``predicting the next word'' even means, as a
probability. Section 3 turns text into token IDs. Section 4 turns IDs into
meaningful, ordered vectors. Section 5 turns those vectors into an actual
predicted word. Nothing here needs attention or Transformers in detail yet.
Those are reviewed, not retaught, starting Session 3 --- if they are not
solid yet, work through the Prerequisites Primer first.

% =============================================================================
\section{Probabilistic language models: predicting the next word, on purpose}
\secsub{A language model is, formally, nothing more than a rule for guessing
the next word from the words said so far.}

Formally, a language model assigns a probability to what word comes next,
given everything written before it:

\begin{equation*}
P(w_t \mid w_1, w_2, \ldots, w_{t-1})
\end{equation*}

\begin{intuition}
This is the whole job, stated precisely. Everything else in this course ---
tokenization, embeddings, attention, the Transformer --- exists only to
compute this one probability well. It runs over and over, one word at a
time.
\end{intuition}

\begin{everyday}
If you have watched every film in a franchise, in order, you can guess what
kind of scene follows a big battle. You do this with decent confidence, not
certainty: usually it is a quiet aftermath, not another battle immediately
after. You did not memorise a rule. You absorbed the pattern from dozens of
examples. A language model does the same thing with text. Having ``read''
huge amounts of it, it estimates how likely each possible next word is,
given everything written so far.
\end{everyday}

\begin{worked}{a tiny count-based language model, by hand}
Suppose your entire training corpus is just four short sentences. They are:
``the dog barked at the mailman'' and ``the dog barked at the cat''. Also:
``the dog barked at the mailman'' again, and ``the dog slept on the mat''.
\begin{steps}
  \item \textbf{Find every exact match for the 5-word context.} The phrase
    ``the dog barked at the'' occurs 3 times in total, across the corpus.
  \item \textbf{Count what follows each time.} It is followed by
    ``mailman'' twice and by ``cat'' once --- 3 continuations observed.
  \item \textbf{Turn counts into probabilities.} So
    $P(\text{mailman}) = 2/3 \approx 0.67$ and
    $P(\text{cat}) = 1/3 \approx 0.33$.
  \item \textbf{Predict.} This tiny model predicts ``mailman''. It is the
    single most likely next word, purely because it appeared more often in
    the small corpus it was shown.
\end{steps}
This is literally what an \vocab{n-gram model} does. It counts how often
each word follows the previous fixed number of words, directly from a
frequency table. No understanding, no meaning --- only counting.
\end{worked}

\housefig[0.6]{figures/fig_ngram_probs.png}{The tiny model's full output: a
probability for every word it has ever seen follow this exact context.}

\begin{table}[htbp]
\centering\small
\begin{tabular}{@{}p{3.1cm}p{4.3cm}p{3.5cm}@{}}
\toprule
\textbf{Approach} & \textbf{How it estimates the probability} & \textbf{Where it breaks} \\
\midrule
N-gram models & Count how often a word follows the previous $N\!-\!1$ words,
  directly from frequency tables & Only ever ``sees'' the last $N\!-\!1$
  words; no sense of meaning \\
Neural language models & A small neural network learns a continuous
  representation of context instead of raw counts & Still a fixed-size
  context window; learns slowly from scratch \\
Transformer-based LLMs & Attention lets the model look arbitrarily far back,
  weighting every earlier word by relevance & The hard problems move to
  cost, data, and reliability (later sessions) \\
\bottomrule
\end{tabular}
\caption*{\footnotesize Three ways to estimate the same probability, in increasing order of capability and cost.}
\end{table}

\begin{watchout}
An n-gram model's blind spot is structural, not a bug you can tune away. If
$N = 4$, it is mathematically incapable of noticing a word five positions
back, no matter how relevant that word is. This is the exact same
fixed-window memory problem that limits a plain RNN, covered in the
Prerequisites Primer.
\end{watchout}

\begin{keytake}
A language model's whole job is estimating $P(w_t \mid w_1, \ldots, w_{t-1})$.
N-grams estimate it by counting. Neural and Transformer models estimate it by
learning a representation --- the rest of this session builds that
representation, piece by piece.
\end{keytake}

% =============================================================================
\section{Tokenization and BPE: text into reusable pieces}
\secsub{A sentence must become a short list of known pieces before any
arithmetic can run on it --- not whole words.}

A model cannot store a separate meaning for every possible word. New
words, typos, and names appear constantly. A fixed whole-word vocabulary
would either explode in size, or fail on anything it has not seen before.
\vocab{Byte Pair Encoding (BPE)} solves this by building a vocabulary of
reusable \emph{sub-word} pieces instead of whole words.

\begin{intuition}
Think of whole-word vocabularies as pre-built plastic statues: one fixed
piece per word, nothing reusable, and no way to build something new. BPE
pieces are LEGO bricks. They are a modest set of small, reusable pieces that
combine to build almost any word, including ones never seen before.
\end{intuition}

\begin{worked}{BPE on the toy corpus \{low$\times$5, lower$\times$2, lowest$\times$2, newer$\times$6, wider$\times$3\}}
\begin{steps}
  \item \textbf{Start from characters.} Every word is split into characters,
    each ending with an end-of-word marker \texttt{\_}. The five words
    become \texttt{l o w \_} ($\times 5$) and \texttt{l o w e r \_}
    ($\times 2$). The next two are \texttt{l o w e s t \_} ($\times 2$) and
    \texttt{n e w e r \_} ($\times 6$). The last is \texttt{w i d e r \_}
    ($\times 3$).
  \item \textbf{Count every adjacent pair across the whole weighted
    corpus.} The pair \texttt{(e, r)} appears in ``lower'' ($\times 2$),
    ``newer'' ($\times 6$), and ``wider'' ($\times 3$) $= 11$ times --- the
    most frequent pair. Merge it: \texttt{e r} $\to$ \texttt{er}.
  \item \textbf{Recount.} Now \texttt{(er, \_)} appears $2+6+3=11$ times
    (every word that just got ``er'' immediately ends in it). Merge it:
    \texttt{er} + \texttt{\_} $\to$ \texttt{er\_}.
  \item \textbf{Recount again.} \texttt{(n, e)} appears in ``newer''
    $\times 6$. But \texttt{(l, o)} appears in ``low'' ($\times 5$),
    ``lower'' ($\times 2$), and ``lowest'' ($\times 2$) --- $9$ times
    total, now the winner. Merge \texttt{l o} $\to$ \texttt{lo}.
\end{steps}
The process continues for a fixed number of merges, or until a target
vocabulary size is reached. Notice what falls out for free. ``Lower'' and
``newer'' end up sharing the learned piece \texttt{er\_}, and ``low'',
``lower'', ``lowest'' end up sharing \texttt{lo}. The algorithm discovers
word-forming patterns purely from frequency --- nobody hand-coded that
``\textit{-er}'' is a comparative suffix.
\end{worked}

\housefig{figures/fig_bpe_merge_counts.png}{The first merge decision: BPE
always merges whichever adjacent pair is currently most frequent across the
whole corpus.}

\begin{everyday}
A rare dog breed name like ``Labradoodle'' is almost never seen whole in
training text. BPE has very likely already learned ``Labra'' (from
Labrador) and ``doodle'' (from Poodle-cross breed names) as separate, common
pieces from other words. The rare whole word is rebuilt from pieces the model
already knows well. The same trick handles typos, brand names, and words in
languages the model has seen very little of.
\end{everyday}

\vocab{Byte-level BPE} (what GPT-style tokenizers actually use) runs this
exact merge process starting from raw UTF-8 \emph{bytes} rather than Unicode
characters. With only 256 possible starting symbols, every byte value, there
is no such thing as an unrepresentable string. Any text in any language or
emoji can always be tokenized, possibly just as more, smaller pieces for rare
scripts.

\begin{waitwhat}
Tokenization is also why LLMs can be strangely bad at some tasks that look
trivial to a human. Counting letters inside a word is a good example. The
model never sees individual characters directly. It only sees whichever
sub-word pieces BPE happened to merge that word into. A letter-counting
question is secretly a question about those pieces, not about reasoning
ability.
\end{waitwhat}

\begin{keytake}
BPE builds a vocabulary of reusable sub-word pieces by repeatedly merging the
most frequent adjacent pair in the training corpus. It handles rare and
unseen words gracefully, at the cost of the model never directly seeing
individual letters.
\end{keytake}

% =============================================================================
\section{Embeddings and positional encoding: meaning, then order}
\secsub{Each token gets a vector that captures its meaning, plus a second
vector that records where it sits in the sentence.}

A token ID on its own is just an arbitrary number --- it carries no sense of
meaning. The next stage looks that ID up in a learned \vocab{embedding
table}. That is a big lookup table that returns a vector of numbers meant
to capture what the token \emph{means}. That meaning is learned entirely
from how the token is used across huge amounts of text.

\begin{intuition}
Two tokens that tend to appear in similar contexts end up with similar
vectors. Nobody labelled them as similar. The model is trained to predict
text well, and giving similar words similar vectors makes that prediction
easier.
\end{intuition}

\begin{everyday}
This is exactly how a streaming service recommends movies. Two films watched
by the same kind of viewer get pulled close together in the service's
internal ``taste space''. This happens even in similar contexts, with
nobody tagging them. A word embedding space works the same way, just for
words instead of films.
\end{everyday}

\begin{worked}{comparing three toy word vectors by hand}
Suppose a (deliberately tiny, two-number) embedding table gives: $\text{dog}
= (0.9, 0.1)$, $\text{cat} = (0.8, 0.2)$, $\text{car} = (0.1, 0.9)$.
\begin{steps}
  \item \textbf{Compute the dot product of dog and cat.}
    $0.9 \times 0.8 + 0.1 \times 0.2 = 0.72 + 0.02 = 0.74$.
  \item \textbf{Compute each vector's length.} $|\text{dog}| =
    \sqrt{0.9^2+0.1^2} = \sqrt{0.82} \approx 0.91$, and $|\text{cat}| =
    \sqrt{0.8^2+0.2^2} = \sqrt{0.68} \approx 0.82$.
  \item \textbf{Divide to get cosine similarity.}
    $0.74 / (0.91 \times 0.82) \approx 0.74/0.75 \approx 0.99$ --- very
    close to 1, meaning ``dog'' and ``cat'' point in almost the same
    direction.
  \item \textbf{Repeat for dog and car.} Dot product $= 0.9 \times 0.1 + 0.1
    \times 0.9 = 0.18$; $|\text{car}| \approx 0.91$; cosine similarity
    $\approx 0.18/0.83 \approx 0.22$ --- far less similar.
\end{steps}
``Dog'' sits close to ``cat'' and far from ``car'' in this toy space. This
holds purely as numbers --- no dictionary definition was involved anywhere
in the calculation.
\end{worked}

\housefig{figures/fig_embedding_space.png}{The toy embedding space from the
worked example: ``dog'' and ``cat'' point in nearly the same direction;
``car'' points elsewhere.}

A separate, lightweight \vocab{embedding model} does exactly this lookup as
its entire job. It is used directly for search and retrieval, not buried
inside a generative LLM. That matters in Session 12's retrieval pipeline,
where running a full LLM just to get a vector would be wasteful.

Meaning alone is not enough, though. ``The dog bit the man'' and ``the man
bit the dog'' use the exact same words and the exact same meaning vectors.
Yet they mean opposite things. A \vocab{positional encoding} vector, one per
position in the sentence, is added on top of each token's meaning vector.
The two are literally summed, element by element, before anything enters
the Transformer stack.

\begin{everyday}
A positional encoding is like the timestamp on a subtitle file. The subtitle
text alone (the meaning) does not say when it appears on screen. The
timestamp (the position) is a separate piece of information, attached to
every single line. The player needs it to show things in the right order.
\end{everyday}

\begin{waitwhat}
Attention itself has no built-in sense of order --- it compares token pairs
regardless of where they sit in the sequence. Without a positional encoding
added in, ``the model'' and ``model the'' would look completely identical to
it. Order is not free. It has to be engineered in, on purpose.
\end{waitwhat}

\begin{keytake}
Each token's final input vector is its meaning (from the embedding table)
plus its position (from positional encoding), added together. Meaning
without position cannot tell word order apart; this sum fixes that.
\end{keytake}

% =============================================================================
\section{The language modelling head: from understanding to one word}
\secsub{One more layer turns the Transformer stack's reasoning into an
actual probability for every word in the vocabulary.}

Every token is first enriched by the Transformer stack's attention and
feed-forward layers, reviewed from Session 3 onward. After that, one more
learned layer takes over. This is the \vocab{language modelling (LM) head}.
It projects the final hidden state into a raw score for every single word
in the vocabulary. That means tens of thousands of numbers. The
\vocab{softmax} function turns those raw scores into probabilities that
sum to exactly 1.

\begin{worked}{softmax on five candidate next words, by hand}
Say the LM head produces raw scores for one prompt: ``The dog movie was
really \_\_\_''. The five candidate scores are great $= 1.80$, good
$=1.35$, slow $=0.60$, long $=0.35$, and boring $=-0.30$.
\begin{steps}
  \item \textbf{Exponentiate every score.} $e^{1.80}\approx 6.05$,
    $e^{1.35}\approx 3.86$, $e^{0.60}\approx 1.82$, $e^{0.35}\approx 1.42$,
    $e^{-0.30}\approx 0.74$.
  \item \textbf{Sum the exponentials.} $6.05+3.86+1.82+1.42+0.74 \approx
    13.89$.
  \item \textbf{Divide each by the sum.} $P(\text{great}) \approx
    6.05/13.89 \approx 0.44$, $P(\text{good})\approx 0.28$,
    $P(\text{slow})\approx 0.13$, $P(\text{long})\approx 0.10$,
    $P(\text{boring})\approx 0.05$. These five add up to $1.00$.
\end{steps}
``Great'' is the single most likely next word here, but it is not certain:
nearly one prediction in three favours ``good'' instead.
\end{worked}

\housefig{figures/fig_next_token_pmf.png}{The worked softmax result: a full
probability distribution over candidate next words, not a single hard
answer.}

A \vocab{decoding strategy} then has to pick one actual word from this
distribution. Always taking the top probability (\vocab{greedy decoding})
is one option, among several covered properly in Session 5. Whichever word
is chosen gets appended to the sequence. The entire five-stage pipeline
from Section 1 runs again to pick the \emph{next} word after that. This
repeat-one-token-at-a-time loop is \vocab{autoregressive generation}. It is
the reason an LLM visibly ``types'' its answer rather than producing the
whole reply instantly.

\begin{watchout}
Always picking the single highest-probability word (greedy decoding) can
trap a model in dull, repetitive text. The single best next word is
sometimes a safe, bland choice, repeated over and over. This is exactly why
Session 5 exists: sampling strategies trade some of that safety for variety.
\end{watchout}

\begin{keytake}
The LM head turns the Transformer's final understanding into a probability
over the entire vocabulary, via softmax. A decoding strategy picks one word
from that distribution. The loop appends it, and the five-stage pipeline
runs again --- one token per pass, from the first word to the last.
\end{keytake}

\noindent\textbf{Where it's used.} Every chat assistant you have used ---
a customer-support bot, a coding assistant, anything --- is running this
exact five-stage loop. It happens underneath its product polish: tokenize,
embed and place, reason through the Transformer stack, score every
candidate word, pick one, repeat.

% =============================================================================
\section*{Self-test}
\addcontentsline{toc}{section}{Self-test}

\noindent\textbf{1.} Why does a 4-gram language model fail to use a relevant
word that appeared six words earlier in the sentence?

\flipanswer{A 4-gram model only ever conditions its prediction on the
previous 3 words ($N-1$ with $N=4$). Anything further back is structurally
invisible to it, regardless of how relevant it is.}

\noindent\textbf{2.} In the BPE worked example, why was \texttt{(e, r)}
merged before \texttt{(l, o)}? ``Low'' is the most frequent whole word, more
frequent than ``newer''.

\flipanswer{BPE counts adjacent \emph{symbol pairs} across the whole
weighted corpus, not whole-word frequency. \texttt{(e, r)} totalled 11
occurrences, from lower, newer, and wider combined. That beat 9 for
\texttt{(l, o)}, so it won the first merge despite ``low'' alone being the
single most frequent word.}

\noindent\textbf{3.} Two tokens end up with nearly identical embedding
vectors. What does that actually mean, mechanically?

\flipanswer{It means the model's training pushed them to behave similarly
for the purpose of predicting surrounding text well. That is not the same
as a human labelling them synonyms. Similar usage in context, not shared
dictionary meaning, is what the vectors capture.}

\noindent\textbf{4.} ``The dog bit the man'' and ``the man bit the dog'' use
identical tokens and identical token embeddings. What makes the model treat
them differently, and why is that necessary?

\flipanswer{Positional encoding, added on top of each token's embedding.
Attention alone compares token pairs with no sense of order. Without a
positional vector, these two sentences would look identical to it.}

\noindent\textbf{5.} The LM head outputs a full probability distribution
over the vocabulary, not a single word. What has to happen next, and why
might always picking the top probability be a bad idea?

\flipanswer{A decoding strategy must pick one word from that distribution.
Always picking the single highest-probability word (greedy decoding) can
produce dull, repetitive text. The safest word is sometimes chosen over and
over.}

% =============================================================================
\section*{Pocket card: Session 1 in one page}
\addcontentsline{toc}{section}{Pocket card}

\begin{keytake}
\textbf{The whole session, compressed.}
\begin{itemize}[leftmargin=1.3em,topsep=4pt,itemsep=4pt]
  \item A language model's entire job is estimating
    $P(w_t \mid w_1, \ldots, w_{t-1})$ --- n-grams count it, neural and
    Transformer models learn a representation instead.
  \item \textbf{Tokenization (BPE)} turns text into reusable sub-word
    pieces by repeatedly merging the most frequent adjacent pair. This
    handles rare and unseen words gracefully.
  \item \textbf{Embeddings} give each token a meaning vector, learned from
    context, not a dictionary. \textbf{Positional encoding}, added on top,
    gives it a place in the sentence. Both are necessary; neither alone is
    enough.
  \item The \textbf{LM head} turns the Transformer's final understanding
    into a probability over the whole vocabulary via softmax. A
    \textbf{decoding strategy} then picks one word.
  \item \textbf{Autoregressive generation}: append the chosen word, run the
    whole five-stage pipeline again, one token per pass.
\end{itemize}
\end{keytake}

% =============================================================================
\section*{Glossary \& symbols}
\addcontentsline{toc}{section}{Glossary \& symbols}

\noindent\textbf{Key terms.}
\begin{description}[leftmargin=9em,style=nextline,font=\normalfont\bfseries\color{navy}]
  \item[language model] a system that assigns a probability to what word
    comes next, given the words said so far.
  \item[n-gram model] a language model that estimates that probability by
    counting how often a word follows the previous fixed number of words.
  \item[token] the smallest unit of text a model actually reads --- usually
    a sub-word piece, not a whole word.
  \item[BPE] Byte Pair Encoding; builds a tokenizer's vocabulary by
    repeatedly merging the most frequent adjacent symbol pair.
  \item[embedding] a learned vector meant to capture a token's meaning,
    based on how it is used in context.
  \item[positional encoding] a vector added to each token's embedding that
    records its position in the sequence.
  \item[Transformer stack] the layers of attention and feed-forward
    computation that process the embedded, positioned sequence.
  \item[LM head] the final learned layer that turns the Transformer's
    output into a raw score for every word in the vocabulary.
  \item[softmax] the function that turns a set of raw scores into
    probabilities that sum to 1.
  \item[decoding strategy] the rule used to pick one actual word from the
    LM head's probability distribution (e.g.\ greedy, sampling).
  \item[autoregressive generation] generating one token per pass, each new
    token appended before the next pass runs.
\end{description}

\medskip
\noindent\textbf{Notation at a glance.}
\begin{center}\small
\begin{tabular}{@{}ll@{}}
\toprule
\textbf{Shorthand} & \textbf{Meaning} \\
\midrule
$w_t$ & the token (word piece) at position $t$ in the sequence \\
$P(w_t \mid w_1,\ldots,w_{t-1})$ & probability of the next token, given every token before it \\
\bottomrule
\end{tabular}
\end{center}

\end{document}
